% =====================================================================
%  Zentrale Präambel für alle Protokolle (Physik F-Praktikum, Synthesechemie)
%  Eingebunden in jedem Protokoll via: \input{../Praeambel.tex}
%  (in den Protokoll-Ordnern, da die Praeambel im Repo-Root liegt)
% =====================================================================

% ---------- Kodierung & Sprache ----------
\AtBeginDocument{\RenewCommandCopy\qty\SI}
\usepackage[utf8]{inputenc} %Um verschiedene Eingabecodierungen zu unterstützen, z.B. UTF-8
\usepackage{fontenc} %Zum Einstellen der Ausgabezeichenkodierung (T1)
\usepackage[ngerman]{babel} %Setzt Sprache auf Deutsch
\usepackage{microtype} %Verbessert die Mikrotypografie (Zeichenabstand, Ränder): ästhetischeres Textbild
% Source Sans Pro Schriftart, eine moderne serifenlose Schriftart.
% Fällt auf lmodern zurück, wenn das Paket nicht installiert ist
% (Installation bei Bedarf: sudo tlmgr install sourcesanspro bzw. apt install texlive-fonts-extra)
\IfFileExists{sourcesanspro.sty}{%
  \usepackage{sourcesanspro}%
}{%
  \usepackage{lmodern}%Vektorisierte Standardschrift, scharf auch im PDF
}
\usepackage{textcomp} %Für zusätzliche Symbole
\usepackage{csquotes} %Typografisch korrekte Anführungszeichen: \enquote{...}
\usepackage{xcolor} %Um Farben (auch Code) darzustellen

% ---------- Seitenlayout ----------
\usepackage{geometry} %Für die Anpassung der Seitenränder und des Seitenlayouts
\geometry{
  top=4cm,
  bottom=3cm,
  left=3cm,
  right=3cm,
  headheight=15pt,
  headsep=30pt,
  footskip=20pt
}

% ---------- Pakete für Physik ----------
\usepackage{amsmath} %Mathematische Erweiterung
\usepackage{amssymb} %Weitere mathematische Zeichen
\usepackage{amsfonts} %Weitere mathematische Schriftarten
\usepackage{amsthm} %Für Beweise und Sätze/Definitionen
\usepackage{mathtools} %Weitergehende Mathe-Werkzeuge
\usepackage{physics} %Zum Einsetzen von physikalischen Einheiten und Operatoren (\abs, \dv, \bra, \ket, ...)
\usepackage{esvect} %Um Vektoren besser darzustellen als mit \vec{}: \vv{} (funktioniert nicht in Captions)
\usepackage{siunitx} %SI-Einheitensystem, z.B. \SI{9.81}{\meter\per\square\second}
\usepackage{bm} %Fette Mathe-Symbole: \bm{\alpha}

% ---------- Pakete für Grafiken und Tabellen ----------
% Für Tabellen
\usepackage{booktabs} %Für schönere Tabellen: \toprule, \midrule, \bottomrule
\usepackage{array} %Für erweiterte Tabellenoptionen
\usepackage{tabularx} %Für Tabellen mit variabler Breite
\usepackage{longtable} %Für mehrseitige Tabellen
\usepackage{multirow} %Für Zellen über mehrere Zeilen

% Für Grafiken
\usepackage{graphicx} %Zur Darstellung von Bildern und Grafiken
\graphicspath{{Bilder/}{../Bilder/}} %Suchpfad für Bilddateien
\usepackage{eso-pic} %Um Bilder im Hintergrund darzustellen
\usepackage{float} %Für exakte Platzierung mit [H]
\usepackage{caption} %Zur Darstellung von Überschriften (Captions)
\usepackage{wrapfig} %Ermöglicht es, Text um Bilder herum zu schreiben
\usepackage{subcaption} %Um Bilder nebeneinander darzustellen mit individuellen Unterschriften
\usepackage[export]{adjustbox} %Um Bilder zu positionieren
\usepackage{tikz} %Zum Selbsterstellen von Grafiken
\usetikzlibrary{shapes, positioning, arrows.meta, calc, graphs, graphs.standard} %Erweiterte TikZ-Bausteine
\usepackage{pgfplots} %Zum Erstellen von Plots direkt in LaTeX
\pgfplotsset{compat=1.18}

% ---------- Pakete für Chemie (nur bei Bedarf aktiv lassen) ----------
\usepackage[version=4]{mhchem} %Chemische Formeln: \ce{H2O}, \ce{2H2 + O2 -> 2H2O}
\usepackage{chemfig} %Strukturformeln: \chemfig{H-C(-[2]OH)(-[4]H)-H}
\usepackage{chemmacros} %Reaktionsumgebungen: \begin{reaction} ... \end{reaction}


% ---------- Pakete für Verweise, Zitate und Quellenangaben ----------
\usepackage[colorlinks=true, allcolors=black]{hyperref} %Für Hyperlinks innerhalb des Dokumentes; alle Links schwarz, damit die kompilierte Datei besser aussieht
\usepackage{cleveref} %Für intelligentere Referenzierungen von Gleichungen, Abbildungen, Tabellen
\usepackage[backend=biber, style=authoryear]{biblatex} %Literaturverwaltung mit biber
\addbibresource{literatur.bib} %Literaturdatenbank im Protokoll-Ordner

% ---------- Pakete zur Struktur ----------
\usepackage{multicol} %Ermöglicht mehrspaltige Textbereiche
\usepackage{fancyhdr} %Für Kopf- und Fußzeilen
\usepackage{titling} %Anpassung von Titel(seiten)
\usepackage{enumitem} %Kompakte, konfigurierbare Listen
\setlist{itemsep=2pt, topsep=4pt} %Engere, aufgeräumte Listen
\usepackage{pdfpages} %Zum Einfügen externer PDFs (z.B. Scans im Anhang)
\usepackage{comment} %Für auskommentierte Notizen: \begin{comment} ... \end{comment}

% ---------- Pakete für Code im Anhang ----------
\usepackage{listings} %Für möglichen Code in den Anhang
\usepackage{mdframed} %Ermöglicht das Erstellen von Rahmen um Text oder andere Inhalte

% Darstellung des Codes im Anhang
\lstset{
    backgroundcolor=\color{white},
    basicstyle=\footnotesize\ttfamily,
    breaklines=true,
    captionpos=b,
    commentstyle=\color{gray},
    frame=single,
    keywordstyle=\color{blue},
    language=Python,
    numbers=left,
    numbersep=5pt,
    numberstyle=\tiny\color{gray},
    showstringspaces=false,
    stringstyle=\color{red},
}

% ---------- Kopf- und Fußzeilen (fancyhdr) ----------
\pagestyle{fancy}
\fancyhf{} %Alle Kopf- und Fußzeilenfelder leeren
\setlength{\headheight}{15.04742pt} %Länge der Kopfzeile festlegen
\addtolength{\topmargin}{-0.04762pt} %Abstand oberer Rand zur Kopfzeile wird verringert
\addtolength{\footskip}{-2.49998pt} %Abstand unterer Rand zur Fußzeile wird verringert
\fancyhead[L]{\rightmark} %Links: aktueller Abschnitt
\fancyhead[R]{\leftmark} %Rechts: aktuelles Kapitel
\fancyfoot[C]{\thepage} %Mitte Fußzeile: Seitenzahl

% ---------- Eigene Befehle ----------
% Unsicherheit: \errval{12.3}{0.4} -> (12.3 ± 0.4)
\newcommand{\errval}[2]{\ensuremath{\left(#1 \pm #2\right)}}
% Mittelwert mit Standardabweichung
\newcommand{\meanval}[2]{\ensuremath{\overline{#1} = #2}}
% Absoluter Fehler einer Größe X
\newcommand{\derr}[1]{\ensuremath{\Delta #1}}
% Gauß'sche Fehlerfortpflanzung als Formelgerüst
\newcommand{\gausserror}{\ensuremath{\Delta y = \sqrt{\sum_i \left(\frac{\partial f}{\partial x_i}\Delta x_i\right)^2}}}
% Häufige Abkürzungen
\newcommand{\ds}{\ensuremath{\displaystyle}}
